Em 29 de outubro de 2026, o grupo conferiu a implementação e os registros de testes das mudanças relacionadas aos relatos de 25 de outubro. A Tabela~\ref{tab:feedback-ajustes} distingue correção técnica, sugestão parcialmente atendida e benefício já existente. Trata-se de uma etapa posterior ao parecer de Thiago; não constitui nova avaliação pelo profissional nem reavaliação com os usuários convidados. Os horários identificam as fontes, não participantes independentes.

\begingroup\small\SingleSpacing\setlength{\parskip}{0pt}\setlength{\LTcapwidth}{\textwidth}
\begin{longtable}{>{\raggedright\arraybackslash}p{3.3cm}>{\raggedright\arraybackslash}p{6.1cm}>{\raggedright\arraybackslash}p{5.2cm}}
\caption{Encaminhamento dos feedbacks e estado da verificação em 29 de outubro.}\label{tab:feedback-ajustes}\\\hline
Origem e observação & Alteração ou benefício preservado & Verificação e limite\\\hline\noalign{\vskip 4pt}\endfirsthead
\hline Origem e observação & Alteração ou benefício preservado & Verificação e limite\\\hline\noalign{\vskip 4pt}\endhead
\hline
\multicolumn{3}{p{\dimexpr\linewidth-2\tabcolsep\relax}}{\footnotesize Fonte: Elaboração do grupo a partir dos relatos de 25/10 e das verificações técnicas de 28 e 29/10/2026. Textos sintetizados, sem atribuição de novas opiniões aos avaliadores.}\\
\endlastfoot
13h31 e 13h40: maior clareza e dois campos para embarque e destino. & Os controles exibem rótulos e o local ou estado da seleção. A troca do papel da busca limpa a confirmação anterior. & Comportamento e disposição em telas móveis e desktop conferidos. São dois controles, não dois campos simultâneos de digitação; compreensão ainda não reavaliada.\\ \noalign{\vskip 8pt}
13h40: melhorar a busca. Replay examinado: sugestão de outra rua. & Na consulta com número e tipo explícito de via, a sugestão precisa corresponder ao nome normalizado da rua. & Caso reconstruído rejeita a rua divergente e aceita a correta. Não demonstra precisão geral do geocodificador nem resolve toda ambiguidade de endereço.\\ \noalign{\vskip 8pt}
13h40: parada e endereço produziram resultados distintos no Sacomã. & Seleção de todas as paradas de ônibus em até 200 m; sem nenhuma, até três mais próximas, preservando o limite de caminhada de 1 km. & Na coordenada examinada, 24 registros candidatos em vez de três. Percursos a pé e horários continuam sujeitos a validação; resolução do episódio histórico não comprovada.\\ \noalign{\vskip 8pt}
13h31: falha ao salvar feedback. & Envio externo liberado; falhas preservam nota e texto ao fechar e reabrir o formulário. Nova tentativa conserva o identificador se o conteúdo não mudar. & Persistência com e sem pesquisa, erros e retentativas verificados. O responsável atribuiu o incidente ao bloqueio de acesso externo; falta vínculo inequívoco com a requisição histórica.\\ \noalign{\vskip 8pt}
Replay do Sacomã: caminhada apresentada com busca incompleta. & Aviso de busca parcial também aparece quando há uma alternativa a pé. & Estado reproduzido no navegador. A advertência informa a limitação; não comprova recuperação de uma alternativa de ônibus.\\ \noalign{\vskip 8pt}
13h40 e 13h47min49s: Olho Vivo integrado; 13h47min49s: ausência de anúncios. & Benefícios existentes preservados, sem atribuir sua criação às correções posteriores. & Integração conferida com respostas simuladas e interface inspecionada. Não assegura disponibilidade permanente da fonte operacional.\\ \noalign{\vskip 4pt}
\end{longtable}
\endgroup

A reconstrução do replay não mostrou um confronto controlado entre parada e endereço: as consultas examinadas utilizaram origens diferentes e plataformas 4 e 5. Seu vínculo com a mensagem das 13h40 não é inequívoco. A ampliação dos acessos responde à concentração de paradas próximas, mas não estabelece a causa única do relato. Os 24 candidatos são registros do catálogo, não necessariamente plataformas físicas distintas. A regra completa de acesso está descrita na Seção~\ref{sec:historico-arquitetura}.

Como complemento definido pelo grupo, o formulário passou a aceitar uma avaliação opcional de uma a cinco estrelas, sem nota pré-selecionada. Com nota, o comentário é opcional; sem nota, permanece necessário. O envio não inicia gravação de navegação. Os registros anteriores continuam sem nota: a nova funcionalidade não produz uma série retrospectiva nem uma medida de satisfação já coletada. O rascunho é preservado ao reabrir o formulário, mas não ao recarregar a página.

A conferência utilizou testes com bases temporárias e respostas controladas: 44 testes de backend, 13 unitários de frontend e 25 verificações de navegador, além de uma regressão automatizada específica do formulário. A regra de acesso às paradas teve 13 testes direcionados e um integrado, incluindo a compatibilidade da partida com o tempo de caminhada. Dois desses testes direcionados retomam casos já incluídos no conjunto de backend, reexecutados após a alteração da seleção de paradas. As contagens identificam conjuntos parcialmente sobrepostos e não devem ser somadas como testes distintos. Esses conjuntos verificam comportamentos determinados e não são amostras de passageiros nem um benchmark de desempenho. As evidências e a identificação das alterações locais sobre a revisão \texttt{34a46f2}, branch \texttt{feat/busca-bidirecional}, estão em \nolinkurl{evidencias-feedbacks-2026-10-29/}, junto ao fonte do relatório. Permanece necessária a reavaliação por tarefas para verificar compreensão dos controles, êxito no planejamento e utilidade das correções.
